% !TEX program = xelatex
\documentclass[UTF8,a4paper,12pt]{ctexart}

\usepackage{geometry}
\geometry{left=2.8cm,right=2.8cm,top=2.6cm,bottom=2.6cm}

\usepackage{amsmath}
\usepackage{amssymb}
\usepackage{booktabs}
\usepackage{graphicx}
\usepackage{float}
\usepackage{caption}
\usepackage{subcaption}
\usepackage{hyperref}
\usepackage{enumitem}
\usepackage{setspace}
\usepackage[numbers]{natbib}

\hypersetup{
  colorlinks=true,
  linkcolor=black,
  citecolor=black,
  urlcolor=blue
}

\setstretch{1.35}
\setlist[itemize]{leftmargin=2em}
\setlist[enumerate]{leftmargin=2em}

\title{\heiti 基于 HMM 与探索性 SCVI 指数的欧亚杓鹬迁徙行为状态识别和关键停歇地评估}
\author{}
\date{\today}

\begin{document}

\maketitle

\begin{abstract}
欧亚杓鹬（\textit{Numenius arquata}）在迁徙过程中依赖一系列停歇地和局部活动地，这些地点对补给、迁徙时机和保育规划具有重要意义。本研究基于 Movebank 数据集 CURLEW\_VLAANDEREN 中 5 只欧亚杓鹬 2020--2025 年 GPS 轨迹，构建从轨迹预处理、行为状态识别、飞行步长重尾特征检验、停歇地或关键活动地识别到探索性保育价值排序的分析流程。原始轨迹共 393,845 条记录，经清洗、1 小时规则化和环境匹配后获得 69,278 条规则化轨迹记录。HMM 输入数据包含 64,385 条记录、2,823 个 burst 和 5 个个体。基于步长和转角拟合 2、3、4 状态 HMM 后，4 状态模型具有最低信息准则（AIC = 176,814，BIC = 177,308），四类状态可解释为停歇、局部活动/觅食、中距离移动和高速飞行。基于停歇和局部活动/觅食状态点，DBSCAN 共识别出 33 个停歇地或关键活动地候选。高速飞行步长具有重尾特征（$\alpha = 2.381$，$x_{\min} = 1.299$），但 power-law 不优于 lognormal 或 exponential 分布，因此不支持严格 L\'evy 飞行解释。停歇时长模型显示，在 2--240 h 迁徙停歇敏感性子集中，\texttt{resource\_z} 和 \texttt{wind\_z} 方向为正但均不显著（\texttt{resource\_z} = 0.038，$p = 0.891$；\texttt{wind\_z} = 0.220，$p = 0.432$；$R^2 = 0.0219$；RMSE = 40.1 h）。基于代理参数的 NDVI 情景模拟显示，当前 NDVI 下平均预测停留时长为 7.54 h，NDVI 下降 5\%、10\% 和 20\% 时分别为 6.01 h、4.69 h 和 2.34 h，NDVI 上升 5\% 时为 9.12 h，但该模拟属于探索性外推。最后，本文构建探索性 SCVI 综合保育价值指数，对 33 个地点进行排序，其中高优先级 7 个、中优先级 10 个、低优先级 16 个，最高 SCVI 为 0.772，中位数为 0.451。总体而言，HMM 行为状态识别是本研究最稳定的结果；L\'evy、最优停止、NDVI 情景模拟和 SCVI 排序均应作为补充性或探索性结果，用于提出后续验证和保育评估假设。在HMM状态识别基础上，进一步构建探索性状态预测与轨迹模拟网页工具，用于展示未来短时状态概率和可能移动包络。
\end{abstract}

\noindent\textbf{关键词：} 欧亚杓鹬；迁徙停歇地；隐马尔可夫模型；重尾步长；最优停止；NDVI；SCVI；保育优先级

\section{引言}

长距离迁徙鸟类在年度生活史中依赖一系列空间上分离的繁殖地、越冬地和迁徙停歇地。停歇地不仅提供休息和能量补给，也影响迁徙时机、飞行距离、后续路线选择以及个体生存和繁殖成功。对于种群下降或面临栖息地压力的迁徙水鸟而言，识别关键停歇地和局部活动热点是制定保护策略的重要基础。

欧亚杓鹬（\textit{Numenius arquata}）是欧洲及东大西洋迁飞通道上具有重要保育意义的迁徙水鸟。其迁徙过程可能受到湿地丧失、土地利用变化、潮滩资源波动和气候背景变化的共同影响。传统地面调查能够提供重要的种群和栖息地信息，但难以连续追踪个体迁徙过程中的行为状态转换。近年来，高时间分辨率 GPS 追踪数据为识别迁徙行为状态和关键活动地点提供了新的可能。

动物移动轨迹通常包含多个隐含行为过程，例如停歇、局部觅食、短距离转移和长距离迁徙飞行。隐马尔可夫模型（Hidden Markov Model, HMM）能够基于步长和转角等观测变量推断不可直接观测的行为状态，因此已广泛用于动物移动生态研究 \citep{michelot2016placeholder}。在迁徙鸟类研究中，HMM 的优势在于可将连续轨迹转化为具有生态解释意义的状态序列，并为后续停歇地识别提供行为基础。

除行为状态识别外，迁徙飞行步长分布也可反映个体移动尺度的异质性。L\'evy flight 或 L\'evy-like 运动常用于讨论动物搜索或移动中的重尾步长特征，但严格 L\'evy 解释需要证明 power-law 分布相对于替代分布具有优势。因此，在迁徙鸟类研究中，对 L\'evy 假说的检验应区分“具有重尾特征”和“严格支持幂律最优模型”。

停歇地停留时长还可能与资源质量、天气条件和离开决策有关。最优停止理论可用于讨论动物何时离开资源斑块，但在实际迁徙系统中，资源代理变量、样本量和地点功能差异常限制其统计检验力度。本研究使用到达时 NDVI 和风支持探索停留时长变化，但将其定位为经验性和探索性模型，而非严格机制检验。

从保育应用角度看，仅识别停歇地仍不足以直接形成保护优先级。关键地点的相对价值可能同时取决于使用强度、停留时间、资源条件和潜在环境脆弱性。因此，本文进一步构建 SCVI（Stopover Conservation Value Index）探索性综合指标，用于将行为状态识别结果转化为候选保育地点排序。本文的核心主线不是并列堆叠多个模型，而是从 HMM 行为状态识别出发，识别停歇地或关键活动地候选，并进一步探索其保育评估价值。

除静态保育评估外，项目还基于状态转移和历史运动行为为参数，预测了未来短期飞行状态和飞行轨迹，可作为迁徙管理和可视化沟通的辅助工具。遗憾的是，由于数据规模限制，我们的预测结果仅具有探索意义，和自然界中的鸟类飞行状态、行为、轨迹还是存在差异的，这也正说明了我们的大自然存在无限可能！

本研究目标包括：第一，基于 GPS 轨迹识别欧亚杓鹬迁徙过程中的主要行为状态；第二，基于低移动强度行为状态识别停歇地或关键活动地候选；第三，评估高速飞行步长是否具有重尾特征，并谨慎检验 L\'evy 假说；第四，探索 NDVI 和风支持是否解释停留时长；第五，构建探索性 SCVI 指数，对候选地点进行综合保育价值排序。

\section{数据与方法}

\subsection{研究对象与数据来源}

本研究以欧亚杓鹬为研究对象，使用 Movebank 数据集 CURLEW\_VLAANDEREN 中的 GPS 追踪记录开展迁徙行为状态识别与关键停歇地评估 \citep{movebankplaceholder}。项目物种代码为 \texttt{curlew}，时间标准统一为 UTC。原始轨迹数据包含个体编号、时间戳、经纬度及 Movebank 附带的传感器字段。根据原始数据摘要，数据共包含 393,845 条记录，覆盖 5 个追踪个体，时间范围为 2020 年 5 月 5 日至 2025 年 10 月 10 日。

环境数据包括 ERA5 再分析温度和风场数据 \citep{era5placeholder}，以及 MODIS NDVI 数据 \citep{modisndviplaceholder}。ERA5 数据用于提取温度、u/v 风分量、风速和顺逆风支持；NDVI 作为资源质量的遥感代理变量。由于 NDVI 不能直接代表欧亚杓鹬实际可利用食物资源，本文将其用于探索性资源分析、情景模拟和综合保育指标构建。

\subsection{轨迹数据清洗与 1 小时规则化}

轨迹预处理包括字段标准化、时间解析、重复记录删除和异常速度过滤。首先识别个体编号、经纬度和时间字段，并将时间解析为 UTC。随后按个体和时间排序，删除同一个体在同一时间的重复记录，并基于 haversine 距离计算相邻点移动速度。项目将最大合理速度阈值设为 35 m/s，超过该阈值的记录被剔除。清洗后仍保留全部 5 个个体。

为满足环境匹配和 HMM 对时间间隔一致性的要求，清洗后的轨迹被规则化至 1 小时时间分辨率。项目中包括基于 1 小时目标间隔和 20 分钟容差的重采样流程，以及按 burst 进行严格 1 小时线性插值的实现。为避免跨越长时间缺口进行不合理插值，严格插值过程中超过 3 小时的时间间隔被视为新的 burst。后续分析使用规则化后的轨迹数据。

\subsection{ERA5 风温数据与 MODIS NDVI 匹配}

ERA5 风温数据按每个规则化 GPS 点的时间和空间位置提取最近网格值。温度由 K 转换为摄氏度；基于 u、v 风分量计算风速；同时根据相邻轨迹点计算移动方位，并将风矢量投影到移动方向上，得到 \texttt{wind\_support}。其中，正值表示顺风支持，负值表示逆风条件。

MODIS NDVI 按个体与时间进行最近邻匹配。最大允许时间差设为 20 天；若最近 NDVI 观测超过该窗口，则该 GPS 点的 NDVI 记为缺失。原始 NDVI 根据数值范围进行缩放，并将异常范围外的值设为缺失。由于 NDVI 匹配存在缺失和个体间质量差异，本文在解释 NDVI 相关结果时保持谨慎。

\subsection{HMM 行为状态识别}

行为状态识别采用 HMM 完成。HMM 输入来自已规则化并匹配环境变量的轨迹数据。按个体检查相邻点时间间隔，小于 0.9 h 或大于 1.1 h 的位置被视为新的 burst，每个 burst 至少包含 10 个点。随后使用经纬度计算步长和转角，并过滤超过 120 km 的步长。标准化温度和风支持作为状态转移协变量。

本研究分别拟合 2、3 和 4 状态 HMM，观测变量为步长和转角，状态转移公式为 \texttt{\textasciitilde{} temp\_z + wind\_support\_z}。模型通过 AIC 和 BIC 进行比较，并以 BIC 最低者作为最佳模型。最佳模型使用 Viterbi 算法解码各轨迹点的最可能状态。状态标签根据平均步长由小到大赋予生态解释，即停歇、局部活动/觅食、中距离移动和高速飞行。需要说明的是，状态标签是基于运动统计特征的模型解释，而不是地面行为观测验证。

\subsection{L\'evy 步长分布分析}

飞行步长分布分析基于 HMM 解码结果。脚本优先选取高速飞行状态的步长作为飞行步长样本；若高速飞行样本不足，则使用飞行状态。为避免极小步长影响拟合，仅保留大于 0.5 km 的飞行步长，并要求样本量不少于 30。

本研究使用 power-law 模型估计幂律指数 alpha 和最小拟合阈值 xmin，并将 power-law 分布分别与 lognormal 和 exponential 分布进行似然比较。若 power-law 不优于替代分布，则不将结果解释为严格 L\'evy flight。该模块在本文中用于描述高速飞行步长的重尾特征。

\subsection{停歇地识别与环境变量提取}

停歇地识别建立在 HMM 行为状态解码结果之上。首先筛选状态标签为停歇和局部活动/觅食的轨迹点，将其作为潜在停歇、觅食或局部活动位置。随后使用 DBSCAN 进行空间聚类，距离度量为 haversine 距离。聚类参数为 eps = 10 km，min\_samples = 3。脚本还将同一个体在同一空间聚类内的连续访问划分为事件，默认允许的最大时间间隔为 12 h，并保留持续时间在配置范围内的访问事件。

每个停歇地或关键活动地候选记录包括个体编号、中心经纬度、到达和离开时间、停留时长、点数、平均 NDVI、到达 NDVI 和风支持等字段。由于空间聚类可能合并跨年度或跨季节重复访问，同一聚类不一定代表单次迁徙停歇事件。因此，本文统一将 DBSCAN 识别的 33 个空间对象称为“停歇地或关键活动地候选”。

\subsection{最优停止/停歇时长模型}

最优停止相关分析用于探索到达时资源质量和风支持是否能够解释停歇地停留时长。由于样本量较小、NDVI 只是资源代理变量，且部分空间聚类具有极长持续时间，本文将该模块定义为探索性停歇时长模型，而不是对最优停止理论的严格机制检验。

模型响应变量为 \texttt{log(duration\_hr)}，解释变量包括到达 NDVI 标准化值和风支持标准化值。主模型使用持续时间为 2--240 h 的迁徙停歇敏感性子集，公式为 \texttt{log(duration\_hr) \textasciitilde{} ndvi\_arrive\_z + wind\_support\_z}。同时，本研究在全部 duration $\geq$ 2 h、2--240 h 迁徙停歇敏感性子集和严格 2--168 h 子集中进行敏感性分析。为衔接后续情景模拟，模型还生成 \texttt{lambda\_proxy} 和 \texttt{Qstar\_proxy} 两个代理参数；这些参数不解释为严格估计的最优停止阈值。

\subsection{NDVI 情景模拟}

NDVI 情景模拟基于停歇时长模型的代理参数，设定当前 NDVI、NDVI 下降 5\%、下降 10\%、下降 20\% 和上升 5\% 五种情景。对每个停歇地和每个情景，按设定倍率调整到达 NDVI，并使用 $\Delta t^{*} = \log(Q_0 / Q^{*}) / \lambda$ 计算预测停留时长；当预测值小于 0 时截断为 0。根据情景 NDVI 与 \texttt{Qstar\_proxy} 的关系划分风险等级。由于基础停歇时长模型解释度较低，该模块被定义为探索性资源情景外推，而不是确定性气候变化预测。

\subsection{SCVI 停歇地保育价值指数}

SCVI 用于对停歇地或关键活动地候选进行探索性保育价值排序。该指数整合四类分项：使用强度、停留时长、资源质量和情景脆弱性。使用强度基于状态点数或停歇点数，经 \texttt{log1p} 转换并归一化；停留时长分项在主分析中设置 240 h 上限，以减弱极端长停留地点对排序的支配；资源质量分项使用到达 NDVI 或平均 NDVI；脆弱性分项来自 NDVI 下降 20\% 情景下相对于当前情景的预测停留时长损失。

主分析权重为使用强度 0.25、停留时长 0.30、资源质量 0.25、脆弱性 0.20。四个分项加权求和得到 SCVI。根据 SCVI 分位数划分保育优先级：前 20\% 为高优先级，50\%--80\% 为中优先级，其余为低优先级。SCVI 在本文中作为探索性综合指标使用，不直接等同于最终保护优先级。

\subsection{探索性状态预测与轨迹模拟网页工具}

在HMM行为状态识别基础上，本研究进一步构建了探索性状态预测与轨迹模拟网页工具，用于展示给定环境和给定位置信息下欧亚杓鹬短期行为状态转移

\section{结果}

\subsection{数据预处理与环境匹配结果}

原始 Movebank 轨迹数据共包含 393,845 条记录，覆盖 5 只欧亚杓鹬个体，时间范围为 2020 年 5 月 5 日至 2025 年 10 月 10 日。原始数据中经纬度和时间字段未出现缺失。经重复记录删除、异常速度过滤和个体记录数筛选后，5 个追踪个体均被保留（表 \ref{tab:preprocess}）。

轨迹经 1 小时规则化后，共获得 69,278 条规则化轨迹记录。规则化检查显示，各个体的中位时间间隔均为 1 h。基于规则化轨迹进一步匹配 ERA5 风温数据和 MODIS NDVI 资源代理变量后，形成带有温度、风速、风支持和 NDVI 的轨迹环境数据。

NDVI 匹配结果显示，在 69,278 条规则化轨迹记录中，17,927 条缺失 NDVI，缺失率为 25.88\%。有效 NDVI 的中位数为 0.485，均值为 0.481，范围为 -0.1975 至 0.8784。个体间 NDVI 匹配质量存在差异（表 \ref{tab:ndviqc}）。

\begin{table}[H]
  \centering
  \caption{原始数据、清洗和规则化结果摘要}
  \label{tab:preprocess}
  \begin{tabular}{ll}
    \toprule
    指标 & 结果 \\
    \midrule
    原始记录数 & 393,845 \\
    个体数 & 5 \\
    时间范围 & 2020-05-05 至 2025-10-10 \\
    1 小时规则化轨迹记录数 & 69,278 \\
    经纬度和时间缺失 & 0 \\
    \bottomrule
  \end{tabular}
\end{table}

\begin{table}[H]
  \centering
  \caption{NDVI 匹配质量摘要}
  \label{tab:ndviqc}
  \begin{tabular}{ll}
    \toprule
    指标 & 结果 \\
    \midrule
    规则化轨迹记录数 & 69,278 \\
    NDVI 缺失数 & 17,927 \\
    NDVI 缺失率 & 25.88\% \\
    NDVI 中位数 & 0.485 \\
    NDVI 均值 & 0.481 \\
    NDVI 范围 & -0.1975 至 0.8784 \\
    \bottomrule
  \end{tabular}
\end{table}

\subsection{HMM 行为状态识别结果}

用于 HMM 建模的输入数据共包含 64,385 条记录，来自 2,823 个 burst，覆盖全部 5 个追踪个体。输入数据中包含 2,823 个步长缺失值和 5,646 个转角缺失值；最大步长为 104.32 km，低于项目设定的 120 km 最大步长阈值（表 \ref{tab:hmminput}）。

2、3 和 4 状态 HMM 的模型比较结果显示，4 状态模型具有最低 AIC 和 BIC。4 状态模型 AIC = 176,814，BIC = 177,308；相比之下，3 状态模型 BIC = 182,805，2 状态模型 BIC = 193,208（表 \ref{tab:hmmcompare}）。后续行为状态结果基于 4 状态 HMM 的 Viterbi 解码。

4 个 HMM 状态在平均步长上呈现梯度，可解释为停歇、局部活动/觅食、中距离移动和高速飞行。停歇状态包含 15,907 条记录，平均步长为 0.0216 km；局部活动/觅食状态包含 33,784 条记录，平均步长为 0.211 km；中距离移动状态包含 8,643 条记录，平均步长为 2.02 km；高速飞行状态包含 405 条记录，平均步长为 46.9 km（表 \ref{tab:states}）。状态轨迹见图 \ref{fig:hmmtracks}。状态标签来自步长和转角特征的模型解释，而非地面行为直接观测。

\begin{table}[H]
  \centering
  \caption{HMM 输入数据摘要}
  \label{tab:hmminput}
  \begin{tabular}{ll}
    \toprule
    指标 & 结果 \\
    \midrule
    HMM 输入记录数 & 64,385 \\
    burst 数量 & 2,823 \\
    个体数 & 5 \\
    步长缺失数 & 2,823 \\
    转角缺失数 & 5,646 \\
    最大步长 & 104.32 km \\
    \bottomrule
  \end{tabular}
\end{table}

\begin{table}[H]
  \centering
  \caption{HMM 候选状态数模型比较}
  \label{tab:hmmcompare}
  \begin{tabular}{ccc}
    \toprule
    状态数 & AIC & BIC \\
    \midrule
    4 & 176,814 & 177,308 \\
    3 & 182,517 & 182,805 \\
    2 & 193,073 & 193,208 \\
    \bottomrule
  \end{tabular}
\end{table}

\begin{table}[H]
  \centering
  \caption{HMM 四状态行为解释和步长摘要}
  \label{tab:states}
  \begin{tabular}{lrr}
    \toprule
    状态解释 & 记录数 & 平均步长（km） \\
    \midrule
    停歇 & 15,907 & 0.0216 \\
    局部活动/觅食 & 33,784 & 0.211 \\
    中距离移动 & 8,643 & 2.02 \\
    高速飞行 & 405 & 46.9 \\
    \bottomrule
  \end{tabular}
\end{table}

\begin{figure}[H]
  \centering
  \includegraphics[width=0.95\textwidth]{\detokenize{../output/figures/08_state_tracks.png}}
  \caption{HMM 行为状态空间轨迹图。}
  \label{fig:hmmtracks}
\end{figure}

\subsection{高速飞行步长分布与 L\'evy 检验}

基于 HMM 解码结果，本研究选取高速飞行状态的步长进行重尾分布和 L\'evy-like 特征检验。power-law 拟合结果显示，幂律指数 $\alpha = 2.381$，$x_{\min} = 1.299$（表 \ref{tab:levy}）。

分布比较结果显示，power-law 与 lognormal 分布比较得到 $R_{\mathrm{lognormal}} = -171.36$，$p_{\mathrm{lognormal}} < 0.001$；与 exponential 分布比较得到 $R_{\mathrm{exponential}} = -80.33$，$p_{\mathrm{exponential}} = 0.005$。两个比较中的 R 值均为负，表明 power-law 并不优于对应替代分布。因此，本研究仅将该结果表述为高速飞行步长具有重尾特征，不将其解释为严格 L\'evy 飞行或纯幂律最优模型。

\begin{table}[H]
  \centering
  \caption{高速飞行步长 power-law 拟合与分布比较}
  \label{tab:levy}
  \begin{tabular}{lr}
    \toprule
    指标 & 结果 \\
    \midrule
    $\alpha$ & 2.381 \\
    $x_{\min}$ & 1.299 \\
    $R_{\mathrm{lognormal}}$ & -171.36 \\
    $p_{\mathrm{lognormal}}$ & $< 0.001$ \\
    $R_{\mathrm{exponential}}$ & -80.33 \\
    $p_{\mathrm{exponential}}$ & 0.005 \\
    \bottomrule
  \end{tabular}
\end{table}

\subsection{停歇地识别与停歇时长模型}

基于 HMM 识别出的停歇和局部活动/觅食状态点，DBSCAN 空间聚类共识别出 33 个停歇地或关键活动地候选。33 个地点均具有有效停留时长，最短停留 2 h，中位数为 15 h，均值为 6,006.82 h，最大值为 45,186 h（表 \ref{tab:duration}；图 \ref{fig:duration}）。

停歇时长模型以到达 NDVI 作为资源质量代理变量，以风支持作为迁徙环境条件变量。主模型使用 2--240 h 的迁徙停歇敏感性子集，样本量 $n = 22$。模型结果显示，\texttt{resource\_z} 的系数为 0.038，$p = 0.891$；\texttt{wind\_z} 的系数为 0.220，$p = 0.432$。两个变量的方向均为正，但均未达到统计显著性。模型整体解释度较低，$R^2 = 0.0219$，RMSE = 40.1 h（表 \ref{tab:mainmodel}；图 \ref{fig:predobs}）。

敏感性分析显示，不同停留时长子集下资源和风支持变量的效应均不显著。全部 duration $\geq$ 2 h 的模型中，NDVI 系数为 -0.272，$p = 0.678$，风支持系数为 0.830，$p = 0.210$；2--240 h 迁徙停歇敏感性子集中，NDVI 系数为 0.038，$p = 0.891$，风支持系数为 0.220，$p = 0.432$；严格 2--168 h 子集中，NDVI 系数为 -0.108，$p = 0.649$，风支持系数为 0.150，$p = 0.528$（表 \ref{tab:sensitivity}）。这些模型均未显示显著效应。

最优停止代理参数估计得到 \texttt{lambda\_proxy} = 0.0253，\texttt{Qstar\_proxy} = 0.478。上述参数仅用于后续探索性情景模拟，不作为显著支持最优停止假说的证据。

\begin{table}[H]
  \centering
  \caption{停歇地或关键活动地候选停留时长摘要}
  \label{tab:duration}
  \begin{tabular}{lr}
    \toprule
    指标 & 结果 \\
    \midrule
    地点数 & 33 \\
    最短停留时长 & 2 h \\
    25\% 分位数 & 5 h \\
    中位数 & 15 h \\
    均值 & 6,006.82 h \\
    75\% 分位数 & 5,523 h \\
    最大值 & 45,186 h \\
    \bottomrule
  \end{tabular}
\end{table}

\begin{figure}[H]
  \centering
  \begin{subfigure}{0.48\textwidth}
    \centering
    \includegraphics[width=\textwidth]{\detokenize{../output/figures/10b_duration_histogram.png}}
    \caption{原始时长分布}
  \end{subfigure}
  \begin{subfigure}{0.48\textwidth}
    \centering
    \includegraphics[width=\textwidth]{\detokenize{../output/figures/10b_log_duration_histogram.png}}
    \caption{log 时长分布}
  \end{subfigure}
  \caption{停歇地或关键活动地候选停留时长分布。}
  \label{fig:duration}
\end{figure}

\begin{table}[H]
  \centering
  \caption{2--240 h 迁徙停歇敏感性子集主模型结果}
  \label{tab:mainmodel}
  \begin{tabular}{lrr}
    \toprule
    指标 & 估计值 & $p$ 值 \\
    \midrule
    \texttt{resource\_z} & 0.038 & 0.891 \\
    \texttt{wind\_z} & 0.220 & 0.432 \\
    \midrule
    样本量 & \multicolumn{2}{c}{22} \\
    $R^2$ & \multicolumn{2}{c}{0.0219} \\
    RMSE & \multicolumn{2}{c}{40.1 h} \\
    \texttt{lambda\_proxy} & \multicolumn{2}{c}{0.0253} \\
    \texttt{Qstar\_proxy} & \multicolumn{2}{c}{0.478} \\
    \bottomrule
  \end{tabular}
\end{table}

\begin{table}[H]
  \centering
  \caption{停歇时长模型敏感性分析}
  \label{tab:sensitivity}
  \begin{tabular}{lrrrrr}
    \toprule
    子集 & $n$ & NDVI 系数 & NDVI $p$ & 风支持系数 & 风支持 $p$ \\
    \midrule
    全部 duration $\geq$ 2 h & 33 & -0.272 & 0.678 & 0.830 & 0.210 \\
    2--240 h 迁徙停歇敏感性子集 & 22 & 0.038 & 0.891 & 0.220 & 0.432 \\
    严格 2--168 h 子集 & 21 & -0.108 & 0.649 & 0.150 & 0.528 \\
    \bottomrule
  \end{tabular}
\end{table}

\begin{figure}[H]
  \centering
  \includegraphics[width=0.75\textwidth]{\detokenize{../output/figures/10_optimal_stopping_pred_obs.png}}
  \caption{停歇时长模型预测值与观测值对比。}
  \label{fig:predobs}
\end{figure}

\subsection{NDVI 资源情景模拟}

在停歇时长模型代理参数基础上，本研究进一步开展 NDVI 资源情景模拟。该模拟属于探索性外推。

当前 NDVI 情景下，模型预测的平均停留时长为 7.54 h。随着 NDVI 下降，预测平均停留时长逐步缩短：NDVI 下降 5\% 时为 6.01 h，下降 10\% 时为 4.69 h，下降 20\% 时为 2.34 h。相反，NDVI 上升 5\% 时，预测平均停留时长增加至 9.12 h（表 \ref{tab:scenario}；图 \ref{fig:scenario}）。

风险分类结果显示，当前情景下高风险、中风险和低风险地点数分别为 5、5 和 12；当 NDVI 下降 20\% 时，高风险地点为 12 个，中风险地点为 5 个，低风险地点为 5 个。上述结果为探索性资源敏感性分析，不解释为确定的生态响应。

\begin{table}[H]
  \centering
  \caption{NDVI 资源情景模拟汇总}
  \label{tab:scenario}
  \begin{tabular}{lrrrr}
    \toprule
    情景 & 平均预测停留时长（h） & 高风险 & 中风险 & 低风险 \\
    \midrule
    当前 NDVI & 7.54 & 5 & 5 & 12 \\
    NDVI 下降 5\% & 6.01 & 6 & 5 & 11 \\
    NDVI 下降 10\% & 4.69 & 10 & 1 & 11 \\
    NDVI 下降 20\% & 2.34 & 12 & 5 & 5 \\
    NDVI 上升 5\% & 9.12 & 4 & 2 & 16 \\
    \bottomrule
  \end{tabular}
\end{table}

\begin{figure}[H]
  \centering
  \includegraphics[width=0.75\textwidth]{\detokenize{../output/figures/11_climate_scenario_projection.png}}
  \caption{不同 NDVI 情景下的平均预测停留时长。该模拟为探索性外推。}
  \label{fig:scenario}
\end{figure}

\subsection{SCVI 停歇地保育价值排序}

本研究将使用强度、停留时长、资源质量和 NDVI 下降情景下的脆弱性整合为 SCVI 停歇地保育价值指数。共有 33 个停歇地或关键活动地候选参与排序，其中高优先级 7 个，中优先级 10 个，低优先级 16 个。SCVI 最高值为 0.772，中位数为 0.451（表 \ref{tab:scvi}）。

排名前 7 的高优先级地点包括 cluster 0、10、18、8、11、2 和 3。其中最高排名地点为 cluster 0，SCVI = 0.772；第二位为 cluster 10，SCVI = 0.726；第三位为 cluster 18，SCVI = 0.696。SCVI 排名见图 \ref{fig:scvibar}，空间分布见图 \ref{fig:scvimap}。

SCVI 在本文中作为探索性综合指标使用，不直接等同于最终保护优先级结论。

\begin{table}[H]
  \centering
  \caption{SCVI 停歇地保育价值排序汇总}
  \label{tab:scvi}
  \begin{tabular}{lr}
    \toprule
    指标 & 结果 \\
    \midrule
    参与排序地点数 & 33 \\
    高优先级地点数 & 7 \\
    中优先级地点数 & 10 \\
    低优先级地点数 & 16 \\
    最高 SCVI & 0.772 \\
    SCVI 中位数 & 0.451 \\
    \bottomrule
  \end{tabular}
\end{table}

\begin{figure}[H]
  \centering
  \includegraphics[width=0.85\textwidth]{\detokenize{../output/figures/12_scvi_ranking_bar.png}}
  \caption{SCVI 排名前列停歇地或关键活动地候选。}
  \label{fig:scvibar}
\end{figure}

\begin{figure}[H]
  \centering
  \includegraphics[width=0.78\textwidth]{\detokenize{../output/figures/12_scvi_stopover_map.png}}
  \caption{停歇地或关键活动地候选的 SCVI 空间分布。}
  \label{fig:scvimap}
\end{figure}

\section{讨论}

\subsection{HMM 四状态模型的行为生态意义}

本研究中，HMM 将欧亚杓鹬轨迹稳定划分为四类运动状态，并呈现出从近乎静止、局部移动、中距离移动到高速飞行的清晰步长梯度。这一结果说明，基于步长和转角的隐状态模型能够捕捉欧亚杓鹬迁徙过程中的主要行为差异，也为后续停歇地识别提供了较可靠的行为基础。与直接基于空间点密度识别停歇地相比，先通过 HMM 排除高速飞行和中距离移动状态，有助于降低将飞行经过区域误判为停歇地的风险。

四状态结构在行为生态上具有合理解释。停歇状态代表低移动强度点，可能对应休息、夜间停留或短暂停顿；局部活动/觅食状态具有较小但连续的移动步长，可能对应觅食、局部搜索或在栖息地内活动；中距离移动状态反映非高速但方向性较强的位移过程，可能对应局地转移、短距离迁移或在不同栖息斑块之间移动；高速飞行状态则更接近迁徙飞行阶段。这样的状态梯度与迁徙鸟类在迁徙路线中交替进行飞行、停留和局部补给的生态过程相符。

不过，HMM 状态仍应被理解为基于运动统计特征的模型分类，而不是对真实行为的直接观测。停歇状态并不必然等同于生理补给，局部活动/觅食状态也不必然证明个体正在觅食。特别是欧亚杓鹬可能在潮滩、农田、湿地和繁殖地等不同环境中表现出相似的低速移动特征。因此，本研究更适宜将 HMM 结果作为识别行为阶段和筛选关键活动地点的基础，而不是直接推断精细行为类型。

\subsection{高速飞行重尾特征与 L\'evy 假说的限制}

高速飞行步长呈现重尾特征，表明欧亚杓鹬迁徙飞行并非由单一尺度的移动距离组成，而是包含少量长距离飞行和更多较短距离位移。这种异质性符合迁徙生态的一般预期：个体可能在部分阶段进行较长距离定向飞行，而在接近停歇地、繁殖地或越冬地时转为较短距离移动。重尾步长也可能反映迁徙路线中风场、地理屏障、停歇地可用性和个体状态共同作用下形成的多尺度移动模式。

然而，本研究结果并不支持将这种重尾特征进一步解释为严格 L\'evy flight。分布比较显示，power-law 模型并不优于 lognormal 和 exponential 等替代分布。因此，虽然步长分布具有长尾，但不能据此推断欧亚杓鹬采用了经典 L\'evy 搜索策略。迁徙鸟类的长距离移动往往受到方向性迁徙、季节性目标、记忆、风场和地形约束影响，并不一定符合随机搜索理论中 L\'evy flight 的前提条件。

此外，L\'evy 分析所使用的是 HMM 识别出的高速飞行状态，而高速飞行状态样本量相对有限。步长分布还可能受 1 小时时间分辨率、轨迹规则化、最大步长过滤和状态分类误差影响。因而，本文更谨慎地将该结果解释为高速飞行距离的重尾特征，而不是迁徙运动机制的严格证明。

\subsection{NDVI 与风支持未显著解释停留时长的可能原因}

停歇时长模型中，NDVI 和风支持的方向在主模型中为正，但均未显著解释停留时长，且模型解释度较低。这一结果可能有多方面原因。首先，样本量限制较明显。尽管轨迹点数量较大，但经停歇地识别和时长筛选后，用于主模型的停歇事件数较少，统计检验能力有限。在样本量较小且停留时长高度右偏的情况下，单个地点或个体的影响可能显著改变模型估计。

其次，NDVI 可能不是欧亚杓鹬食物资源的直接代理。欧亚杓鹬的食物资源与潮滩底栖动物、土壤无脊椎动物、湿地水文条件、潮汐暴露时间和土地利用方式等因素有关。NDVI 主要反映植被绿度或地表生产力，在农田或草地环境中可能与潜在觅食条件有一定关系，但在潮滩、泥滩或开阔湿地中，其与可利用食物资源之间的关系可能较弱甚至不一致。因此，NDVI 与停留时长之间未出现显著关系，并不必然意味着资源不重要，而可能说明当前资源代理变量不够贴合欧亚杓鹬的真实补给资源。

第三，风支持对停留时长的影响可能并非简单线性关系。顺风条件可能促进个体更快离开停歇地，也可能降低飞行成本并改变后续迁徙决策；逆风条件可能延长等待时间，但也可能与天气系统、降水或气压变化共同作用。仅使用停歇期间平均风支持或到达附近的风支持，可能无法充分捕捉个体离开决策所依赖的未来风场窗口。此外，风效应可能与迁徙方向、季节、个体状态和地理位置交互，而当前模型未能纳入这些复杂结构。

第四，停歇地定义中可能混有长期栖息地、越冬地、繁殖相关地点或跨年度重复利用地点。部分空间聚类持续时间极长，说明它们不完全代表单次迁徙过程中的补给停歇事件。即使模型中使用 2--240 h 子集进行主分析，停歇地识别和环境汇总仍可能受到空间聚类定义的影响。若不同生态功能的地点被放入同一模型，资源和风支持对停留时长的真实关系可能被稀释。

因此，最优停止模型结果应被解读为当前数据和代理变量下尚未发现显著解释关系，而不是资源质量和风场不影响欧亚杓鹬停留决策。更严格的检验需要更多停歇事件、更精确的资源变量、明确区分迁徙停歇与长期栖息地，并可能需要分层模型或个体随机效应来处理个体差异。

\subsection{NDVI 情景模拟对迁徙补给窗口的启示}

NDVI 情景模拟显示，在当前代理模型设定下，NDVI 下降会压缩模型预测的停留时长，并使更多地点被标记为接近或低于资源阈值。这一结果的主要意义不在于预测未来具体停留时长，而在于提供一种思考迁徙补给窗口脆弱性的框架。若停歇地资源质量下降，个体可用于补给的有效窗口可能缩短，进而影响迁徙节律、后续飞行准备和迁徙路线选择。

从保育角度看，这类模拟有助于识别潜在敏感地点。即使当前模型并未显著证明 NDVI 决定停留时长，情景模拟仍能作为假设生成工具：哪些地点在资源下降假设下更容易接近阈值，哪些地点可能需要优先监测资源变化，哪些地点的保护价值可能不仅取决于当前使用强度，也取决于未来资源稳定性。

不过，本研究的 NDVI 情景模拟必须被视为探索性外推。其输入参数来自解释度较低的停歇时长模型，资源变量为 NDVI 代理指标，且情景设置为固定比例扰动，而非基于真实未来气候模型或生态过程模型。模拟结果还没有纳入潮汐、水文、人类干扰、保护地管理和食物资源动态等关键因素。因此，该部分结果更适合用于提出后续监测和情景分析问题，而不应作为确定性预测或因果推断。

\subsection{SCVI 对关键停歇地识别的应用价值}

SCVI 将使用强度、停留时长、资源质量和情景脆弱性整合到同一指标中，为从行为状态识别结果走向保育应用提供了一个可操作框架。与单独依据停留时长或点位密度排序相比，SCVI 同时考虑地点被使用的强度、个体停留的持续性、资源代理条件以及在资源下降情景下的潜在敏感性。这有助于将轨迹数据转化为候选保护地点清单，服务于后续实地调查、湿地管理和保护优先级讨论。

本研究中，高 SCVI 地点通常具有较高使用强度、较长持续时间或较高资源质量。对于保护实践而言，这些地点可作为后续核查对象，例如与保护地边界、湿地名录、土地利用类型、人类干扰强度和潮滩资源调查结果叠加分析。SCVI 还可用于比较不同地点的相对关注价值，从而帮助有限监测资源优先投向更可能具有生态意义的区域。

但 SCVI 本身仍是探索性综合指标，而非最终保护决策工具。首先，SCVI 对停留时长和使用强度敏感，可能受到极端长停留地点的影响。虽然项目中对用于 SCVI 的时长设置了上限，但长期栖息地、越冬地或繁殖相关地点仍可能获得较高分值。其次，资源质量和脆弱性分项依赖 NDVI 及其情景模拟，而这些变量本身存在代理不确定性。再次，SCVI 权重具有一定主观性，不同权重设置可能改变地点排序。

因此，SCVI 的合理用途是形成候选重点地点，而非最终保护优先级。在正式管理应用前，应进一步验证高分地点的生态功能，区分迁徙停歇地、繁殖地、越冬地和长期栖息地，并将 SCVI 结果与独立生态调查、保护地现状和威胁评估结合。

\subsection{研究局限与未来改进}

本研究的首要限制是样本个体数较少。尽管轨迹记录量较大，追踪个体仅有 5 只，且个体间追踪时长和空间覆盖差异明显。因此，结果主要代表这些被追踪个体的行为模式，不能简单外推至整个欧亚杓鹬种群或所有迁飞通道。未来研究应整合更多个体、更多年份和不同繁殖或越冬来源的轨迹数据，以检验 HMM 状态结构和关键地点排序是否稳定。

第二，停歇地事件数较少，且停留时长分布高度不均衡。用于最优停止主模型的事件数仅为有限子集，降低了统计检验能力。未来可通过更严格的事件拆分方法，将空间热点按迁徙季节、访问年份和连续停留过程重新划分，避免将多年重复访问合并为单一事件。同时，应明确区分迁徙停歇地、长期栖息地、越冬地和繁殖相关地点。

第三，资源和环境变量仍较粗略。NDVI 不是欧亚杓鹬食物资源的直接代理，ERA5 风场也难以完全代表个体实际飞行高度和局地气象条件。后续研究可引入潮汐暴露时间、底栖动物资源、生境类型、土地利用、人类干扰、水位和保护地管理数据，以更贴近欧亚杓鹬实际补给生态。若可获得加速度、活动传感器或实地观测数据，也可用于验证 HMM 状态标签。

第四，最优停止模型解释度较低，当前结果不足以支持明确的停留决策机制。未来可尝试分层模型、个体随机效应、季节分组、迁徙阶段分组和非线性模型，检验资源、风场和个体状态之间的交互影响。也可将离开时风场、未来 24--72 小时风窗口、前一段飞行距离和后续飞行距离纳入模型，以更接近迁徙停留决策的真实过程。

第五，气候/NDVI 情景模拟目前是探索性外推。后续应将其与真实气候情景、土地利用变化、湿地水文模型或资源动态模型结合，而不是仅依赖固定比例 NDVI 扰动。同时，应对代理参数不确定性进行传播分析，评估不同参数假设下风险地点是否保持一致。

第六，SCVI 可能受极端长停留地点和权重设定影响。虽然时长截尾和敏感性方案可以减弱这一问题，但仍需要系统评估不同权重、不同停歇地定义和不同资源变量对排序结果的影响。未来可将 SCVI 与独立保护价值指标进行交叉验证，例如保护地覆盖程度、已知重要湿地、重复访问年数、多个个体共享程度和人类威胁强度。

总体而言，本研究提供了一条从 HMM 行为状态识别到关键活动地点筛选，再到探索性保育价值评估的分析路径。其较强结论在于 HMM 能够稳定识别欧亚杓鹬轨迹中的主要行为状态，并支持关键活动热点的初步筛选；其较弱和探索性部分则包括 L\'evy 假说、最优停止模型、NDVI 情景外推和 SCVI 排序。未来工作应围绕事件定义、资源测量、样本扩展和独立验证进一步强化这一框架。

\section{结论}

\begin{enumerate}
  \item 本研究基于 5 只欧亚杓鹬 2020--2025 年 GPS 轨迹，建立了从行为状态识别到停歇地或关键活动地候选筛选，再到探索性保育价值排序的分析流程。
  \item 4 状态 HMM 是最稳定的主结果，可将轨迹解释为停歇、局部活动/觅食、中距离移动和高速飞行四类行为状态。
  \item 基于 HMM 低移动强度状态和 DBSCAN 识别出 33 个停歇地或关键活动地候选，但其中可能混有长期栖息地、越冬地或繁殖相关地点。
  \item 高速飞行步长具有重尾特征，但不支持严格 L\'evy 飞行；NDVI 和风支持对停留时长方向在主模型中为正但不显著，不能作为强机制证据。
  \item NDVI 情景模拟和 SCVI 排序可作为探索性保育评估工具，用于提出后续监测和验证的候选地点，而不应直接等同于最终保护决策。
\end{enumerate}

\begin{thebibliography}{99}

\bibitem[Michelot et al.(2016)]{michelot2016placeholder}
Michelot et al., 2016. Hidden Markov models for animal movement analysis. 文献信息待补充。

\bibitem[Movebank]{movebankplaceholder}
Movebank. Movebank animal tracking database and CURLEW\_VLAANDEREN dataset. 文献信息待补充。

\bibitem[ERA5]{era5placeholder}
ERA5. ERA5 reanalysis climate data. 文献信息待补充。

\bibitem[MODIS NDVI]{modisndviplaceholder}
MODIS NDVI. MODIS vegetation index product documentation. 文献信息待补充。

\bibitem[Powerlaw]{powerlawplaceholder}
Powerlaw. Power-law distribution fitting and comparison methods. 文献信息待补充。

\bibitem[DBSCAN]{dbscanplaceholder}
DBSCAN. Density-based spatial clustering method. 文献信息待补充。

\bibitem[Optimal stopping theory]{optimalstoppingplaceholder}
Optimal stopping theory. Optimal stopping and patch departure theory. 文献信息待补充。

\end{thebibliography}

\appendix

\section{附录说明}

附录 A：项目数据处理流程和脚本对应表，列出轨迹清洗、规则化、环境匹配、HMM、L\'evy、停歇地识别、停歇时长模型、NDVI 情景模拟和 SCVI 计算的脚本、输入和输出。

附录 B：HMM 模型补充结果，包括 2、3、4 状态模型 AIC/BIC、状态步长和转角摘要，以及后续可补充的状态转移矩阵。

附录 C：NDVI 和 ERA5 匹配质控，包括 NDVI 缺失率、个体间匹配差异和风支持计算说明。

附录 D：L\'evy 分析详细结果，包括 alpha、xmin、与 lognormal 和 exponential 的分布比较，并说明 power-law 未优于替代分布。

附录 E：33 个停歇地或关键活动地候选完整清单，包括个体编号、中心经纬度、到达和离开时间、停留时长、NDVI、风支持和点数。

附录 F：停歇时长模型敏感性分析，包括全部 duration $\geq$ 2 h、2--240 h 迁徙停歇敏感性子集和严格 2--168 h 子集的完整系数表。

附录 G：NDVI 情景模拟详细结果，包括每个地点在当前、NDVI 下降和 NDVI 上升情景下的预测停留时长、变化量和风险分类。

附录 H：SCVI 权重和排名稳定性分析，包括不同时长上限和不同权重设定下的排序变化。

\end{document}
